% p.31 — الباب السادس عشر في آتى يؤتي (ناقص في أفعل) — تتمة
\mspage{31}
% Layout: six columns, each a continuous stream of paradigms (14 forms each,
% read top to bottom, then the next column to the left). A paradigm begins
% with a boxed heading cell holding the title and, below it in the same box,
% the uncontracted base form; both are given here as two consecutive rows.
% In col. 5 the boxed sections (منفي بالنون … نهي غائب بالنون) show only the
% base form and the first form. Internal box rules are not reproduced.
% Hamza on initial alif is written as printed (bare alif + ḍamma: اُوْتِ).
\begin{sarf}{|c|c|c|c|c|c|}
لَمْ تُؤْتِيَا & لَمْ اُوْتَ & لَنْ يُؤْتَيَا & لَتُؤْتِيَنَّ & لَتُؤْتَيْنَانِّ & \thead{أمر حاضر بالنون}\\
لَمْ تُؤْتُوْا & لَمْ نُؤْتَ & لَنْ يُؤْتَوْا & لَتُؤْتِيَانِّ & لَاُوْتَيَنَّ & اَأْتِيَنَّ\\
لَمْ تُؤْتِيْ & \thead{مضارع بلن} & لَنْ تُؤْتَى & لَتُؤْتُنَّ & لَنُؤْتَيَنَّ & آتِيَنَّ\\
لَمْ تُؤْتِيَا & لَنْ يُؤْتِيَ & لَنْ تُؤْتَيَا & لَتُؤْتِنَّ & \thead{منفي بالنون} & آتِيَانِّ\\
لَمْ تُؤْتِيْنَ & لَنْ يُؤْتِيَ & لَنْ يُؤْتَيْنَ & لَتُؤْتِيَانِّ & لَا يُؤْتِيَنَّ & آتُنَّ\\
لَمْ اُوْتِ & لَنْ يُؤْتِيَا & لَنْ تُؤْتَى & لَتُؤْتِيْنَانِّ & لَا يُؤْتِيَنَّ & آتِنَّ\\
لَمْ نُؤْتِ & لَنْ يُؤْتُوْا & لَنْ تُؤْتَيَا & لَاُوْتِيَنَّ & \thead{مجهول بالنون} & آتِيَانِّ\\
\thead{مجهول بلم} & لَنْ تُؤْتِيَ & لَنْ تُؤْتَوْا & لَنُؤْتِيَنَّ & لَا يُؤْتَيَنَّ & آتِيْنَانِّ\\
لَمْ يُؤْتَى & لَنْ تُؤْتِيَا & لَنْ تُؤْتَيْ & \thead{مجهول بالنون} & لَا يُؤْتَيَنَّ & \thead{نهي حاضر بالنون}\\
لَمْ يُؤْتَ & لَنْ يُؤْتِيْنَ & لَنْ تُؤْتَيَا & لَيُؤْتَيَنَّ & \thead{أمر غائب بالنون} & لَا تُؤْتِيَنَّ\\
لَمْ يُؤْتَيَا & لَنْ تُؤْتِيَ & لَنْ تُؤْتَيْنَ & لَيُؤْتَيَنَّ & لِيُؤْتِيَنَّ & لَا تُؤْتِيَنَّ\\
لَمْ يُؤْتَوْا & لَنْ تُؤْتِيَا & لَنْ اُوْتَى & لَيُؤْتَيَانِّ & لِيُؤْتِيَنَّ & لَا تُؤْتِيَانِّ\\
لَمْ تُؤْتَ & لَنْ تُؤْتُوْا & لَنْ نُؤْتَى & لَيُؤْتَوُنَّ & \thead{مجهول بالنون} & لَا تُؤْتُنَّ\\
لَمْ تُؤْتَيَا & لَنْ تُؤْتِيْ & \thead{مضارع بالنون} & لَتُؤْتَيَنَّ & لِيُؤْتَيَنَّ & لَا تُؤْتِنَّ\\
لَمْ يُؤْتَيْنَ & لَنْ تُؤْتِيَا & لَيُؤْتِيَنَّ & لَتُؤْتَيَانِّ & لِيُؤْتَيَنَّ & لَا تُؤْتِيَانِّ\\
لَمْ تُؤْتَ & لَنْ تُؤْتِيْنَ & لَيُؤْتِيَنَّ & لَيُؤْتَيْنَانِّ & \thead{نهي غائب بالنون} & لَا تُؤْتِيْنَانِّ\\
لَمْ تُؤْتَيَا & لَنْ اُوْتِيَ & لَيُؤْتِيَانِّ & لَتُؤْتَيَنَّ & لَا يُؤْتِيَنَّ & \thead{تمرين}\\
لَمْ تُؤْتَوْا & لَنْ نُؤْتِيَ & لَيُؤْتُنَّ & لَتُؤْتَيَانِّ & لَا يُؤْتِيَنَّ & آوٰى\\
لَمْ تُؤْتَيْ & \thead{مجهول بلن} & لَتُؤْتِيَنَّ & لَتُؤْتَوُنَّ & \thead{مجهول بالنون} & آلٰى\\
لَمْ تُؤْتَيَا & لَنْ يُؤْتَى & لَتُؤْتِيَانِّ & لَتُؤْتَيِنَّ & لَا يُؤْتَيَنَّ & آنٰى\\
لَمْ تُؤْتَيْنَ & لَنْ يُؤْتَى & لَيُؤْتِيْنَانِّ & لَتُؤْتَيَانِّ & لَا يُؤْتَيَنَّ & آذٰى\\
\end{sarf}
